% §4 Methodology
% Status: Revised 2026-03-23 — full rewrite reflecting E1–E7 actual protocols.

\section{Methodology}
\label{sec:methodology}

This section describes the experimental design, workload selection, abstraction
implementations, measurement protocols, and profiling methodology used to
systematically characterize abstraction effectiveness. A central feature of this
methodology is its empirical grounding: several protocol decisions reported here
were not assumed \emph{a priori} but were discovered and validated during
experimental execution. Where the protocol departs from common practice, we
explain the empirical motivation. The complete experimental artifact---benchmark
source code, build scripts, raw CSV data, analysis notebooks, and profiling
traces---will be released on Zenodo with a permanent DOI.

% -----------------------------------------------------------------------
\subsection{Experimental Design: Three-Dimensional Matrix}
\label{subsec:3d-matrix}
% -----------------------------------------------------------------------

We structure our experiments as a three-dimensional matrix crossing workload
type, abstraction layer, and target architecture, ensuring systematic coverage
of the abstraction effectiveness space.

\begin{table}[h]
\centering
\caption{Three-dimensional experimental matrix.}
\label{tab:3d-matrix}
\begin{tabular}{p{2.5cm}p{5.5cm}p{5.5cm}}
\hline
\textbf{Dimension} & \textbf{Categories} & \textbf{Justification} \\
\hline
Workloads (7)
  & Regular, Semi-irregular, Irregular
  & Full spectrum of computational predictability \\[4pt]
Abstractions (4)
  & Kokkos, RAJA, SYCL, Julia (+ native baseline)
  & Dominant portability approaches in production codes~\cite{Evans2022} \\[4pt]
Architectures (2)
  & NVIDIA RTX~5060 (Blackwell), AMD MI300X (CDNA3)
  & Two GPU architectures with distinct memory technologies and execution
    models \\
\hline
\end{tabular}
\end{table}

The full matrix yields over \textbf{190 measured configurations} across 3--4
problem sizes per workload and 30 timed runs per configuration. Not all cells
are populated: SYCL is unavailable on the RTX~5060 for most kernels (no
AdaptiveCpp CUDA backend at time of testing), and Numba is excluded entirely
due to compute capability~12.0 / PTX~9.2 incompatibility
(Section~\ref{subsec:abstraction}).

The matrix design enables controlled comparison: by holding two dimensions
constant while varying the third, we isolate the contribution of each factor
to observed performance differences.

% -----------------------------------------------------------------------
\subsection{Target Hardware Platforms}
\label{subsec:hw-method}
% -----------------------------------------------------------------------

Hardware specifications and architectural differences are described in
Section~\ref{subsec:hardware} (Table~\ref{tab:hardware}). The NVIDIA RTX~5060
serves as a Blackwell architecture portability target; the AMD MI300X serves as
the primary HPC platform. The two platforms differ substantially in memory
technology (GDDR7 vs.\ HBM3), measured bandwidth (288 vs.\ 5300\,GB/s), and
thread execution model (warp-32 vs.\ wavefront-64). This platform class
asymmetry is intentional: the study evaluates abstraction portability
\emph{behavior} across architectures, not cross-platform throughput
comparability. PPC values measure efficiency relative to native on each
platform independently (Equation~\ref{eq:efficiency}).

% -----------------------------------------------------------------------
\subsection{Workload Selection}
\label{subsec:workloads}
% -----------------------------------------------------------------------

We select seven kernels spanning the full spectrum of computational regularity,
because prior work suggests that abstraction effectiveness is strongly
correlated with workload structure~\cite{Godoy2023, Evans2022}. Rather than
implementing kernels from scratch, we leverage established benchmark suites to
minimise implementation bias.

\subsubsection{Regular Workloads}

\textbf{E1 --- STREAM Triad} (BabelStream~\cite{Deakin2019}):
$A[i] = B[i] + s \cdot C[i]$. Perfectly regular, bandwidth-bound
($\sim$0.25\,FLOP/byte), embarrassingly parallel. Establishes the performance
ceiling for memory-bound portability: if an abstraction cannot achieve
near-native performance here, it cannot achieve it anywhere.

\textbf{E2 --- Dense Matrix Multiplication (DGEMM)} (RAJAPerf Suite, Julia
microbenchmarks~\cite{Godoy2023}): Compute-bound ($\sim$16--64\,FLOP/byte
depending on blocking), sensitive to memory layout and shared-memory tiling
quality. Tests whether abstractions can express optimal algorithms
(Section~\ref{sec:taxonomy}, pattern P004).

\textbf{E3 --- 3D Stencil} (7-point Jacobi): Structured grid, memory-bound
(AI $\approx 0.203$\,FLOP/byte). Tests multi-dimensional indexing abstractions
and tiling policy overhead (Section~\ref{sec:taxonomy}, pattern P006).

\subsubsection{Semi-Irregular Workloads}

\textbf{E4 --- Sparse Matrix-Vector Multiplication (SpMV)} (CSR format):
Irregular memory access driven by sparsity pattern. We test three matrix classes
--- structured (2D/3D Laplacian), power-law (Pareto $\gamma = 2.5$), and random
(Erd\H{o}s--R\'enyi) --- to isolate the effect of load imbalance on abstraction
overhead (Section~\ref{sec:taxonomy}, pattern P007).

\textbf{E5 --- Sparse Triangular Solve (SpTRSV)} (level-set parallel):
Inherent sequential dependencies require level-set scheduling, producing many
short kernel launches per solve. Tests dispatch amplification overhead
(Section~\ref{sec:taxonomy}, pattern P008).

\subsubsection{Irregular Workloads}

\textbf{E6 --- Breadth-First Search (BFS)} (level-synchronous, GAP Benchmark):
Extreme load imbalance, control-flow divergence, and atomic operations.
Represents the stress test for portability frameworks.

\textbf{E7 --- N-Body Force Calculation} (Lennard-Jones, CSR neighbor list):
Irregular neighbor lists with cutoff radius. High arithmetic intensity but
irregular memory access. Tests dispatch amortization across problem sizes.

\begin{table}[h]
\centering
\caption{Problem sizes per experiment. Each size is chosen to expose a distinct
performance regime: small (launch-overhead-dominated), medium (balanced), and
large (bandwidth- or compute-saturated).}
\label{tab:problem-sizes}
\begin{tabular}{lllll}
\hline
\textbf{Experiment} & \textbf{Small} & \textbf{Medium} & \textbf{Large}
& \textbf{XLarge} \\
\hline
E1 STREAM   & $2^{20}$ (1M)    & $2^{26}$ (67M)    & $2^{28}$ (268M)  & --- \\
E2 DGEMM    & $N = 1024$       & $N = 4096$        & $N = 8192$       & --- \\
E3 Stencil  & $32^3$           & $128^3$           & $256^3$          & $512^3$ \\
E4 SpMV     & structured 1K    & power-law         & scale 20--22     & --- \\
E5 SpTRSV   & structured       & random            & circuit          & --- \\
E6 BFS      & RMAT scale~20    & road network      & Erd\H{o}s--R\'enyi & --- \\
E7 N-Body   & 32K atoms        & 256K atoms        & 2M atoms         & --- \\
\hline
\end{tabular}
\end{table}

% -----------------------------------------------------------------------
\subsection{Abstraction Layers and Implementations}
\label{subsec:abstractions-method}
% -----------------------------------------------------------------------

Table~\ref{tab:abstractions} summarises the four abstraction layers evaluated
and their backends on each platform.

\begin{table}[h]
\centering
\caption{Abstraction layers evaluated in this study.}
\label{tab:abstractions}
\begin{tabular}{lllll}
\hline
\textbf{Layer} & \textbf{Version} & \textbf{NVIDIA Backend}
& \textbf{AMD Backend} & \textbf{Paradigm} \\
\hline
Native CUDA/HIP & 12.8 / ROCm 6.x & ---  & ---  & Explicit GPU programming \\
Kokkos          & 4.x              & CUDA & HIP  & C++ performance portability \\
RAJA            & 2024.x           & CUDA & HIP  & Loop-level policy abstraction \\
SYCL (AdaptiveCpp) & 0.9.4+       & ---  & HIP (HSA) & Open-standard heterogeneous \\
Julia           & 1.10+            & CUDA.jl & AMDGPU.jl & High-level JIT language \\
\hline
\end{tabular}
\end{table}

\textbf{Native CUDA/HIP:} Hand-optimised implementations from benchmark suites
serve as the performance ceiling on each platform. Best practices are enforced:
coalesced memory access, occupancy tuning, and shared memory usage where
applicable.

\textbf{Kokkos~\cite{Carter2014}:} C++ performance portability library with
\texttt{View}s and \texttt{ExecutionSpace}s. \texttt{TeamPolicy} exposes
scoped shared-memory allocation (equivalent to CUDA \texttt{\_\_shared\_\_}),
enabling tiled algorithms where the workload requires it. Widely adopted in ECP
codes ($\sim$30\% per Evans et al.~\cite{Evans2022}).

\textbf{RAJA~\cite{Hornung2014}:} Abstracts loop-level execution policies via
templated \texttt{forall} and reduction operations, while leaving memory
management explicit. RAJA provides no portable API for scoped shared-memory
allocation; kernels requiring tiling must either use vendor extensions or forgo
tiling entirely, incurring an algorithmic memory traffic penalty on
compute-bound workloads (pattern P004).

\textbf{SYCL~\cite{Reyes2020} via AdaptiveCpp:} Open Khronos standard for
heterogeneous programming. \texttt{local\_accessor} provides portable
shared-memory access, making SYCL expressive for tiling-dependent kernels.
On AMD hardware, AdaptiveCpp routes launches through the HSA runtime,
introducing a per-launch dispatch overhead ($\approx$19\,µs) that becomes
significant for fine-grained, level-set workloads (pattern P008). SYCL is
unavailable on our RTX~5060 target via AdaptiveCpp; SYCL results are therefore
AMD-only for most kernels.

\textbf{Julia~\cite{Bezanson2017}:} High-level language with JIT compilation
via LLVM. CUDA.jl and AMDGPU.jl generate device code from high-level array
operations and custom GPU kernels. Julia arrays use column-major storage by
default, which creates architecture-specific coalescing effects (pattern P005).

\textbf{Numba (excluded):} Numba's CUDA JIT backend does not support Blackwell
compute capability~12.0 (PTX~9.2 mismatch). Excluded from all experiments.

% -----------------------------------------------------------------------
\subsection{Measurement Protocol}
\label{subsec:measurement}
% -----------------------------------------------------------------------

\subsubsection{Warmup Protocol}
\label{subsubsec:warmup}

The warmup protocol is one of the most consequential methodological decisions
in this study, revised twice based on empirical evidence.

\paragraph{Standard practice.}
The standard warmup recommendation in the literature is 10 pre-timed
iterations, intended to amortise JIT compilation overhead and warm
caches~\cite{Godoy2023}.

\paragraph{E1 revision: memory-bound protocol.}
Analysis of raw per-run throughput data from E1 (STREAM Triad) revealed that
10 warmup iterations are insufficient on the RTX~5060. The GDDR7 memory
controller undergoes an irreversible clock transition (Section~\ref{subsec:thermal})
requiring approximately 40 sustained high-bandwidth iterations to trigger.
With only 10 warmup runs, timed measurements begin before the GPU has reached
its stable operating frequency. We adopt a \textbf{minimum of 50 pre-timed
warmup iterations} for all memory-bound experiments on the RTX~5060.

\paragraph{E2 revision: compute-bound protocol.}
Sustained FP64 load on the RTX~5060 causes progressive SM frequency reduction
that is not resolved by a fixed iteration count. For compute-bound kernels,
warmup continues until per-iteration throughput variance over a sliding window
of 10 runs falls below 2\% coefficient of variation. This typically requires
60--80 iterations.

\paragraph{MI300X protocol.}
The MI300X uses HBM3 with no dynamic clock behavior. Standard 10-iteration
warmup is sufficient. The extended protocol is applied conservatively for
consistency across all platforms.

\subsubsection{Statistical Rigor}

\begin{itemize}
    \item \textbf{Repetitions:} 30 timed executions per configuration,
    following the warmup protocol.
    \item \textbf{Reporting:} Median execution time and interquartile range
    (IQR). We never report only the mean: GPU kernels exhibit outliers from
    OS jitter and DVFS that inflate mean values.
    \item \textbf{Outlier handling:} Runs with \texttt{hw\_state\_verified = 0}
    (deviating $> 15\%$ from session median, or captured during a thermal
    transition) are excluded before computing statistics.
    \item \textbf{System state:} CPU frequency governor locked to performance
    mode. GPU SM clocks locked via \texttt{nvidia-smi --lock-gpu-clocks}
    (NVIDIA) or \texttt{rocm-smi --setperflevel high} (AMD) before each
    session. Memory clocks on the RTX~5060 cannot be locked directly; the
    warmup protocol compensates.
\end{itemize}

\subsubsection{Baseline Definition}

The native baseline for PPC computation is defined as:

\begin{quote}
\emph{The median throughput of the native CUDA or HIP implementation, measured
in the same hardware thermal regime as all abstraction runs in the same
experimental session, with \texttt{hw\_state\_verified = 1} for all included
runs.}
\end{quote}

\noindent This definition departs from the common practice of using a single
native measurement as a session-independent reference. The departure is
motivated by E1: when native and abstraction runs are measured in different
thermal regimes on the RTX~5060, efficiency ratios reflect hardware state
differences rather than abstraction overhead, producing spurious variations of
up to $\sim$29\%. By requiring thermal regime consistency, we ensure that all
efficiency values are causally attributable to the abstraction layer.

% -----------------------------------------------------------------------
\subsection{Four-Level Measurement Hierarchy}
\label{subsec:hierarchy}
% -----------------------------------------------------------------------

Every configuration proceeds through a four-level measurement hierarchy,
with deeper levels triggered by performance anomalies at the preceding level.
Figure~\ref{fig:methodology-pipeline} illustrates the pipeline.

\begin{figure}[t]
\centering
\includegraphics[width=\columnwidth]{figures/methodology2.png}
\caption{Four-level measurement hierarchy. Every configuration enters at
Level~1; deeper levels are triggered when abstraction efficiency falls below
$0.85\times$ native. Known API limitations (P004) bypass the profiling queue.}
\label{fig:methodology-pipeline}
\end{figure}

\paragraph{Level~1 --- Performance Portability Coefficients.}
The PPC (Equation~\ref{eq:ppc}) is computed for every configuration. Efficiency
is always relative to measured native performance, not theoretical peak. PPC
$\geq 0.80$ denotes excellent portability; $0.60$--$0.80$ acceptable;
$< 0.60$ poor (Section~\ref{subsec:ppc}).

\paragraph{Level~2 --- Overhead Attribution.}
Any configuration where abstraction achieves $< 0.85 \times$ native
performance is automatically flagged for deep profiling. Overhead is decomposed
into: (i) kernel launch latency and dispatch overhead, (ii) memory transfer
overhead (explicit copies, page faults, framework-induced extra copies), and
(iii) resource utilisation gaps (occupancy, achieved bandwidth, L2 cache hit
rate).

An intentional exception applies: configurations where low efficiency is an
expected consequence of API limitations (e.g., RAJA and Julia DGEMM, where
the API cannot express shared-memory tiling) are excluded from the profiling
queue---their low efficiency is a taxonomy finding (pattern P004), not an
implementation defect.

Profiling tools: Nsight Systems (\texttt{nsys}) for timeline analysis and
Nsight Compute (\texttt{ncu}) for kernel-level hardware counters on NVIDIA;
\texttt{rocprof} for hardware counters and Omniperf for memory hierarchy
analysis on AMD.

\paragraph{Level~3 --- Root Cause Analysis.}
Performance gaps identified at Level~2 are classified into one of four root
cause categories:

\begin{enumerate}
    \item \textbf{Compiler backend failure:} suboptimal code generation (missed
    loop unrolling, poor vectorisation). Diagnosed by inspecting generated
    PTX (NVIDIA) or GCN (AMD) assembly side-by-side with the native baseline.
    \item \textbf{Runtime coordination overhead:} extra synchronisation,
    excessive dispatch frequency, or dispatch amplification in level-set
    workloads.
    \item \textbf{Memory model mismatch:} abstraction forces non-optimal memory
    layout or access pattern (e.g., non-coalesced access due to column-major
    storage).
    \item \textbf{API limitation:} the abstraction's programming model cannot
    express a required optimisation (e.g., scoped shared-memory allocation for
    tiling).
\end{enumerate}

\noindent These four categories form the basis of the empirical taxonomy of
eight failure and success patterns developed in Section~\ref{sec:taxonomy}.

\paragraph{Level~4 --- Roofline Characterisation.}
For every configuration, we compute hardware utilisation:

\begin{equation}
    u = \frac{\text{Throughput}_\text{measured}}
             {\min(P_\text{peak},\; \text{AI} \times B_\text{peak})}
    \label{eq:roofline}
\end{equation}

\noindent where $P_\text{peak}$ is peak compute throughput, AI is arithmetic
intensity, and $B_\text{peak}$ is peak memory bandwidth. A native utilisation
gate is enforced: native implementations must achieve $u \geq 0.80$ on the
large problem size before abstraction experiments begin. This verifies that we
are measuring abstraction overhead, not hardware or software setup defects.

% -----------------------------------------------------------------------
\subsection{Thermal Regime Characterisation}
\label{subsec:thermal}
% -----------------------------------------------------------------------

A critical methodological finding emerged during E1 that is absent from prior
performance portability studies: the RTX~5060's consumer-class GDDR7 memory
controller exhibits \emph{irreversible intra-session clock transitions} that
constitute a first-class experimental confound.

\subsubsection{E1 Discovery: Dynamic Memory Clock Transitions}

Under sustained memory-bandwidth load, the RTX~5060 undergoes a one-way
transition from a lower memory clock state ($\sim$9{,}001\,MHz, yielding
$\sim$271\,GB/s) to a higher boost state ($\sim$12{,}001\,MHz, yielding
$\sim$350\,GB/s). This transition is:

\begin{enumerate}
    \item \textbf{Irreversible within a session:} once triggered, the GPU does
    not return to the lower clock state until the driver session is fully reset.
    \item \textbf{Triggered by sustained load:} approximately 40 consecutive
    high-bandwidth iterations are required, regardless of GPU temperature.
    \item \textbf{Abstraction-dependent in timing:} JIT-compiled abstractions
    (Julia) delay the transition by 5--10 iterations relative to ahead-of-time
    compiled frameworks, because JIT compilation latency reduces effective
    memory pressure during early warmup.
\end{enumerate}

\noindent If different abstractions are benchmarked without explicit clock
control, they may be measured in different thermal regimes, producing spurious
performance differences of up to $\sim$29\% ($350 / 271 \approx 1.29$) that are
entirely attributable to hardware state rather than abstraction overhead.

\subsubsection{E2 Discovery: Compute-Bound Thermal Contamination}

E2 (DGEMM) revealed a second class of thermal confound: sustained FP64 load
causes progressive SM frequency reduction, producing a monotonically decreasing
throughput curve within a session. Fresh-state profiling of two flagged
configurations revealed:

\begin{itemize}
    \item \texttt{raja\_naive}/medium: raw $\eta = 0.24$ $\rightarrow$
    thermal-corrected $\eta = 0.66$ ($2.76\times$ thermal ratio). The residual
    gap is a genuine API limitation (pattern P004), not thermal.
    \item \texttt{julia\_naive}/medium: raw $\eta = 0.80$ $\rightarrow$
    corrected $\eta \approx 1.0$. The sub-unity efficiency was entirely
    thermal; Julia's layout advantage (pattern P005) is confirmed
    size-independent.
\end{itemize}

\subsubsection{Control Mechanisms}

\begin{itemize}
    \item \textbf{Hardware state flag (\texttt{hw\_state\_verified}):} Each run
    is tagged with a binary flag. Runs deviating $> 15\%$ from session median
    are flagged as \texttt{0} and excluded from all statistics.
    \item \textbf{SM clock locking:} \texttt{nvidia-smi --lock-gpu-clocks}
    eliminates boost-clock variability in the compute domain.
    \item \textbf{Adaptive warmup:} The CV$< 2\%$ protocol
    (Section~\ref{subsubsec:warmup}) ensures thermal stability before timed
    measurement begins.
    \item \textbf{Same-session baselines:} Native and abstraction runs are
    always measured in the same session under the same thermal conditions.
\end{itemize}

\noindent These thermal confounds are specific to the RTX~5060's GDDR7
memory controller and dynamic power management. The MI300X (HBM3, fixed clocks)
exhibits no such behaviour and requires no special thermal protocol.

% -----------------------------------------------------------------------
\subsection{Data Collection and Automation}
\label{subsec:automation}
% -----------------------------------------------------------------------

All experiments are driven by a Python automation framework
(\texttt{run\_experiments.py}) that iterates over the full experimental matrix.
Per-kernel shell scripts handle binary discovery, warmup execution, timed
measurement, and CSV output. The automation always executes native baselines
first, then abstraction variants in a fixed order, with configurable cooldown
periods between runs.

Each timed run produces one row in a per-abstraction CSV file with the
following schema: \texttt{timestamp}, \texttt{experiment\_id}, \texttt{kernel},
\texttt{abstraction}, \texttt{platform}, \texttt{problem\_size},
\texttt{run\_id}, \texttt{execution\_time\_ms}, \texttt{throughput},
\texttt{hw\_state\_verified}. Raw CSVs are append-only and never edited
post-collection. A validation script (\texttt{validate\_schema.py}) enforces
required columns, enum constraints, positivity, and uniqueness before any
analysis proceeds.

PPC values, profiling queue flags, and roofline utilisation are computed
post-hoc by dedicated analysis scripts (\texttt{compute\_ppc.py},
\texttt{compute\_roofline.py}) that filter on
\texttt{hw\_state\_verified = 1} before computing any statistics. The Git
commit hash of the benchmark suite is recorded at run time; the complete
experimental artifact will be released with a permanent DOI.
